\begin{afterword}
\summaryhead

The post recalls the author's view, from ``Bind Yourself to Reality,'' that a rationalist should pour hope and care into a lawful universe with no magic in it. It borrows the old trick of listing things one is grateful for, and proposes a variant: list the abilities you have that would seem amazing if they were magic, or if only a few people had them.

Binocular vision becomes a magic second eye that sees into the third dimension, the ``Mystic Eyes of Depth Perception.'' Speech becomes Vibratory Telepathy, writing becomes Psychometric Tracery, the hands and tool-making become Multidimensional Kinesis, and sight becomes the Eye. The last and longest item is the Ultimate Power, intelligence: an inner echo of the universe that lets its bearer steer toward any future, overcomes every lesser power, cannot be taken away without taking the bearer, and is held by its bearers to ground moral worth, by an argument the post calls self-serving. Its bearers rule all other life.

\responsehead

The post is a list of ordinary abilities described as fantasy powers, and it should be read as one. Under the fantasy language, most of the descriptions are accurate.

The proposal in its opening is the only claim that asks to be believed, and it is modest: why not make such a list? The trick it varies has some support. In experiments by Emmons and McCullough, people who listed things they were grateful for showed higher well-being on several, though not all, measures. The variant, admiring one's abilities by imagining them as magic, is not tested, and the post does not claim it has been.

The last item takes up about two fifths of the post. It restates ``The Power of Intelligence'' in the language of fantasy: intelligence as a model of the world used to search for paths to chosen futures, the power that overcomes all others. The post then gives, as the view of those who hold the Ultimate Power, the claim that intelligence is part of what makes a creature an end in itself. It calls their defense of it, that no one without the power can understand its importance, ``a suspiciously self-serving argument,'' but does not say whether the claim itself is true.

In short: a list of ordinary abilities recast as magic, which suggests a variant of the gratitude list that has not been tested, and whose longest item restates the author's earlier claims about intelligence.
\end{afterword}
